
qu'il résulte de 7.16.2 que \(\Gr_\bullet^{\mathrm{num}}(X_{\bar t})\) est isomorphe à un quotient d'un sous-groupe de \(\Gr_\bullet^{\mathrm{num}}(X_{\bar s})\) (\(f\) propre et lisse); en particulier, son rang est encore majoré par le rang de ce dernier.

\textbf{7.17.4.} Lorsque \(S\) est régulier en \(s\) et que \(t\) est maximal, on peut construire \(S'\) en faisant éclater \(s\) dans le localisé de \(S\) en \(s\), et en localisant au point générique de la fibre de \(s\) dans l'éclaté. On trouve alors pour \(k(s')\) une \emph{extension pure} de \(k(s)\). Or il n'est pas difficile de montrer que si \(k'\) est une extension pure d'un corps \(k\), alors pour tout schéma noethérien \(T\) sur \(k\), l'homomorphisme canonique \(K_\bullet(T)\to K_\bullet(T\otimes_k k')\) est un isomorphisme, en représentant \(K'\) comme limite inductive des anneaux affines d'ouverts affines de schémas de la forme \(\Spec k[T_1,\ldots,T_n]\), et utilisant (IX 1.6). Donc on conclut de ceci que (7.17.1) peut être considéré comme un homomorphisme canonique
\begin{equation}\tag{7.17.4.1}
  K_\bullet(X_t)\longrightarrow K_\bullet(X_s).
\end{equation}
Mais on notera qu'on ne peut en déduire un homomorphisme \(K_\bullet(X_{\bar t})\to K_\bullet(X_{\bar s})\), car passant à une extension finie \(K_i\) de \(K=k(t)\), et normalisant \(S\) dans \(K_i\), l'hypothèse de régularité sur \(S\) sera perdue.

Un autre façon de définir un homomorphisme (7.17.4.1) consiste à se ramener au cas où \(S\) est local de point fermé \(s\) (ce qui est trivial), et à considérer une suite croissante de sous-schémas fermés \emph{réguliers} \(S_i\) de \(S\), de point générique \(t_i\), tels que \(\dim S_i=i\) (\(0\leq i\leq n=\dim S\)). Comme pour tout \(i\), \(0\leq i\leq n-1\), \(\underline{O}_{S_{i+1},t_i}\) est un anneau de valuation discrète de corps résiduel \(k(t_i)\), corps des fractions \(k(t_{i+1})\), on dispose d'un homomorphisme de spécialisation
\[
  K_\bullet(X_{t_{i+1}})\longrightarrow K_\bullet(X_{t_i}),
\]
